% !PRINTED 426
\weberpage{426}
%<<CONT
is substituted, one finds
\[
D = 5^5(108a^5c-135a^4b^2+90a^2bc^2 - 320ab^3c+256b^5+c^4),
\]
so that
\weq{16}{5^5\Delta = a^2D}
results, and from (15) one obtains for $\lambda$ the expression
\[
\lambda = \frac{11a^3b - ac^2 + 2b^2c + \frac{1}{25}a\sqrt{\frac{1}{5}D}}{2(a^4+abc-b^3)}
\]
in which the square root may take both signs.

One sees that, besides the square root of the discriminant, which is rationally expressible in terms of the roots, $\sqrt5$ also occurs in it.

Once $\lambda$ has been computed, one obtains from (11) the ratio $z_1:\mu^2$, and from (12) the last unknown $z$, represented \gesperrt{rationally} in terms of $\lambda$.

\begin{center}
\S.\,113.

\gesperrt{Resolvents of the Sixth Degree.}
\end{center}

To find the resolvents of degree $6$ of the icosahedral equation, we start from one of the dihedral groups $D_5$ contained in the icosahedral group (\S.\,111, 2.). The representation of the icosahedral group [\S.\,58, (20)] gives us at once the decomposition into the subsidiary groups; if
\weq{1}{Q = \Theta^r,\ \Theta^r\psi,\ r=0,1,2,3,4}
is the dihedral group from which we start, then we obtain the full icosahedral group
\weq{2}{P = Q + Q\chi + Q\chi\Theta + Q\chi\Theta^2 + Q\chi\Theta^3 + Q\chi\Theta^4.}
If now $U_\infty$ is a function which remains unchanged by the substitutions of $Q$ and which is carried by
\[
\chi,\ \chi\Theta,\ \chi\Theta^2,\ \chi\Theta^3,\ \chi\Theta^4
\]
into
\[
U_0,\ U_1,\ U_2,\ U_3,\ U_4
\]
then the $U_\infty, U_0, U_1, U_2, U_3, U_4$ are the roots of a resolvent of degree $6$. (On the notation $U_\infty$ cf. \S.\,65.)
